\NSPage{133}%
On the pulse \NSi{NS-F-P133-001}{XE^2=X_{\mathrm p} e_b^2 f^2 e^{-2\lambda y}}. Changing to
\NSi{NS-F-P133-002}{\zeta_b=\lambda y}, and using \eqref{eq:A.14}, \eqref{eq:A.15}, and
\eqref{eq:A.18}, gives the exact form
\NSFormula{NS-F-P133-003}{display}{A.19}
\begin{equation}
 \frac{\lambda S(\infty)}{X_{\mathrm p} e_b^2 f^2}
 =\operatorname{Amp}^2K_b-\frac{1-e^{-26}}{4}+\mathcal E(\operatorname{Amp},\eta),
 \qquad
 K_b=\int_0^{13}e^{-2\zeta_b}R_0(\zeta_b)^2\,d\zeta_b,
 \tag{A.19}\label{eq:A.19}
\end{equation}
where on \NSi{NS-F-P133-004}{[.9,1.2]\times[-1,1]},
\NSFormula{NS-F-P133-005}{display}{none}
\[
 \lVert\mathcal E\rVert_{C^1_\eta}
 \leq C_{\mathrm{pre}}\lambda(1+\log(1/\lambda)).
\]
All terms contributed by the outer profile \NSi{NS-F-P133-006}{E} in this estimate were
chosen independently of \NSi{NS-F-P133-007}{\operatorname{Amp}}; the exponentially small affine
end bumps account for its remaining dependence. Since
\NSi{NS-F-P133-008}{R_0\leq\zeta_b},
\NSi{NS-F-P133-009}{K_b\leq\int_0^\infty\zeta_b^2e^{-2\zeta_b}\,d\zeta_b=1/4}. On
\NSi{NS-F-P133-010}{.02\leq\zeta_b\leq9} one has
\NSi{NS-F-P133-011}{R_0\geq\zeta_b-.02}, and integration gives
\NSi{NS-F-P133-012}{K_b>.20}. The principal expression in \eqref{eq:A.19} is below
\NSi{NS-F-P133-013}{-.047} at \NSi{NS-F-P133-014}{.9} and above
\NSi{NS-F-P133-015}{.038} at \NSi{NS-F-P133-016}{1.2}, and its derivative in the
amplitude is at least \NSi{NS-F-P133-017}{.36} on the bracket. For small
\NSi{NS-F-P133-018}{\lambda} the full expression has a unique root there. Its linearization
is nonsingular, so the root is smooth and
\NSFormula{NS-F-P133-019}{display}{A.20}
\begin{equation}
 |\partial_\eta\operatorname{Amp}|
 \leq C_{\mathrm{pre}}\lambda(1+\log(1/\lambda)).
 \tag{A.20}\label{eq:A.20}
\end{equation}
Both equalities of \eqref{eq:A.8} now hold. The remaining task is to verify that these moment
corrections preserve the stress cone. Only finitely many orders of
\NSi{NS-F-P133-020}{\eta}-derivatives must satisfy prescribed smallness bounds below. If
bounds on additional finite derivative orders are prescribed, the differentiated integrals and
Lemma~\ref{lem:A.2} provide them before the parameters are fixed. Derivatives of every other
fixed order have finite bounds after the parameters are chosen.

\subsection{Fixing the pressure datum for the axis profile.}\label{sec:A.4}
The azimuthal profile \NSi{NS-F-P133-021}{E} now determines the axis pressure
\NSi{NS-F-P133-022}{\Pi_0} by the normalization that pressure vanish at radial infinity, as in
\eqref{eq:4.25}. The analytic axis construction of Proposition~\ref{prop:B.2} requires this
datum to extend holomorphically near \NSi{NS-F-P133-023}{[-1,1]}, with the sign and size
bounds stated in Lemma~\ref{lem:A.5}. In this subsection we verify these properties before
attaching the axis, and show that the datum is independent of the later choice of the radial
scale \NSi{NS-F-P133-024}{X_R}. For the pressure integrals in this subsection, we use the
global logarithmic coordinate \NSi{NS-F-P133-025}{y=\log(X/X_R)} on the complete reference
profile.

\begin{lemma}\label{lem:A.5}
Let \NSi{NS-F-P133-026}{E_{\mathrm{id,sched}}} be the reference inner profile \eqref{eq:A.7}
followed by the outer profile \NSi{NS-F-P133-027}{E} specified in Section~\ref{sec:A.2}, with
the azimuthal corrections that preserve the total pressure increment omitted and all other
transition intervals unchanged. The function
\NSFormula{NS-F-P133-028}{display}{A.21}
\begin{equation}
 \Pi_0(\eta)=-\frac12\int_{-\infty}^{\infty}
 E_{\mathrm{id,sched}}(y,\eta)^2\,dy
 \tag{A.21}\label{eq:A.21}
\end{equation}
is independent of \NSi{NS-F-P133-029}{X_R} and analytic on a complex neighborhood of
\NSi{NS-F-P133-030}{[-1,1]}. It is even, \NSi{NS-F-P133-031}{\Pi_0'(\eta)} has the sign
of \NSi{NS-F-P133-032}{\eta} for \NSi{NS-F-P133-033}{\eta\neq0}, and
\NSFormula{NS-F-P133-034}{display}{A.22}
\begin{equation}
 \Pi_0(\eta)\leq-\frac52P_*^2f(\eta)^2.
 \tag{A.22}\label{eq:A.22}
\end{equation}
For the complete reference radial profile, integration of
\NSi{NS-F-P133-035}{\Pi_X=E^2/(2X)} from this datum gives
\NSFormula{NS-F-P133-036}{display}{A.23}
\begin{equation}
 \Pi(X,\eta)=-\frac12\int_{\log(X/X_R)}^\infty E(v,\eta)^2\,dv.
 \tag{A.23}\label{eq:A.23}
\end{equation}
\end{lemma}
\begin{proof}
Without the angular bumps, every part of \NSi{NS-F-P133-037}{E} has the form
\NSi{NS-F-P133-038}{c(y)f(\eta)^{\vartheta(y)}}, where \NSi{NS-F-P133-039}{c} and the
transition lengths are independent of \NSi{NS-F-P133-040}{\eta} and
\NSi{NS-F-P133-041}{0\leq\vartheta\leq1}. The reference inner profile, axial transition,
intermediate power-law interval, and pulse have \NSi{NS-F-P133-042}{\vartheta=1}; profile
interpolation decreases it to zero; exterior transition and tail have
\NSi{NS-F-P133-043}{\vartheta=0}. Choose one simply connected
\NSPage{134}%
complex neighborhood of the real interval avoiding the poles and zeros of
\NSi{NS-F-P134-001}{f}, and fix an analytic logarithm there. Then
\NSi{NS-F-P134-002}{f^\vartheta=\exp(\vartheta\log f)} is holomorphic with a bound uniform
in \NSi{NS-F-P134-003}{0\leq\vartheta\leq1}. The squared integrand is bounded by
\NSi{NS-F-P134-004}{Ce^{y/5}} as \NSi{NS-F-P134-005}{y\to-\infty} and by
\NSi{NS-F-P134-006}{Ce^{-(1+2h)y}} as \NSi{NS-F-P134-007}{y\to+\infty}, and the
intermediate interval is finite. Uniform integration of holomorphic functions proves
analyticity.

For real \NSi{NS-F-P134-008}{\eta}, differentiating each factor gives
\NSi{NS-F-P134-009}{\partial_\eta f^{2\vartheta}=-2\vartheta J_0'f^{2\vartheta}}. Thus
every contribution to \NSi{NS-F-P134-010}{\Pi_0'} has the sign of
\NSi{NS-F-P134-011}{\eta}, and the reference inner profile makes it strict. Its contribution
is exactly \NSi{NS-F-P134-012}{-(5/2)P_*^2f^2}, proving \eqref{eq:A.22}. The construction
in the logarithmic radial coordinate contains no \NSi{NS-F-P134-013}{X_R}, proving its
independence. Finally, \eqref{eq:A.11} says that reinstating the angular bumps changes the
total pressure increment by zero. Its total increment is therefore
\NSi{NS-F-P134-014}{-\Pi_0}, which proves \eqref{eq:A.23}.
\end{proof}

We can now use this pressure datum in the axis construction. Subsequent profile changes must
preserve the total pressure increment to keep it fixed. Matching all five radial integrals also
preserves the radial velocity and the functions \NSi{NS-F-P134-015}{Q_s,N_s} beyond the
correction interval. If an edit has zero total integral of
\NSi{NS-F-P134-016}{(E_{\mathrm{new}}^2-E_{\mathrm{old}}^2)/(2X)}, the forward pressure is
unchanged before and after its support, though it can change inside that support. Accordingly,
a correction supported strictly to the right of \NSi{NS-F-P134-017}{X} that preserves the
total pressure increment may be omitted in the backward integral \eqref{eq:A.23}. Equality of
all five radial moment functions at a radius where the fields rejoin has the stronger consequence
of Lemma~\ref{lem:4.4}: it restores the radial velocity and the functions
\NSi{NS-F-P134-018}{Q_s,N_s} on the entire later interval, including their parameter
derivatives. The axis construction will use the already fixed analytic function
\NSi{NS-F-P134-019}{\Pi_0} and integrate pressure forward. Later smooth, possibly
nonanalytic, edits of \NSi{NS-F-P134-020}{E} cannot change that datum; their total pressure
increment is either zero or is restored by the correction matching all five radial moments. The
heat replacement preserves the required convergent moment differences while changing the
exterior profile.

\subsection{Uniform admissible stress cone inequalities on the outer interval.}\label{sec:A.5}
The moment identities and the axis pressure datum are now fixed. It remains to verify the
stress cone inequalities on the outer intervals specified in Proposition~\ref{prop:A.4}. We
estimate the ratios entering the sufficient cone test of Lemma~\ref{lem:4.5} before choosing
the large radial scale \NSi{NS-F-P134-021}{X_R}; this final choice then gives uniform strict
inequalities on the required finite interval.

Write \NSi{NS-F-P134-022}{w=N_s/(EQ_s)} wherever
\NSi{NS-F-P134-023}{Q_s>0}. The sufficient test of Lemma~\ref{lem:4.5} is
\NSFormula{NS-F-P134-024}{display}{A.24}
\begin{equation}
 a-b_sw>0,\qquad 2b_sw+b_s^2/a+(a-2)w^2<2.
 \tag{A.24}\label{eq:A.24}
\end{equation}
We establish uniform positive lower bounds for \NSi{NS-F-P134-025}{a-b_sw} and
\NSi{NS-F-P134-026}{2-2b_sw-b_s^2/a-(a-2)w^2}, with
\NSi{NS-F-P134-027}{a,b_s,w} in fixed compact ranges and
\NSi{NS-F-P134-028}{Q_s>0}. Increasing \NSi{NS-F-P134-029}{X_R} afterward then makes
\NSi{NS-F-P134-030}{p_{s,1}=XQ_s/L} large enough on the designated finite logarithmic
interval. Where \NSi{NS-F-P134-031}{v_s\leq2}, it is enough to have
\NSi{NS-F-P134-032}{P_{\mathrm c}=p_{s,1}(1-b_sw/a)>2}.

The reference inner profile and axial transition. On \eqref{eq:A.7}, direct substitution into
the equation for \NSi{NS-F-P134-033}{Q_s} gives
\NSFormula{NS-F-P134-034}{display}{A.25}
\begin{equation}
 Q_s=\frac{(3/5)(4L-1)-h(1-8\eta^2)+(D+4d)\eta J_0'}{8/5}\geq c>0.
 \tag{A.25}\label{eq:A.25}
\end{equation}
There \NSi{NS-F-P134-035}{b_s=0} and \NSi{NS-F-P134-036}{a=4/5}. In the equation for
\NSi{NS-F-P134-037}{N_s}, split the source into its geometric and pressure parts. At the start
of the first transition their values are respectively \NSi{NS-F-P134-038}{O(|\eta|)} and
\NSi{NS-F-P134-039}{O(P_*^2|\eta|)}; one parameter derivative removes the factor
\NSi{NS-F-P134-040}{|\eta|}. These estimates follow by integrating the constant ideal
geometric source and the bounded pressure source against the
\NSPage{135}%
kernel in the integral representation of \NSi{NS-F-P135-001}{N_s}. Evenness of the pressure
and the ideal shapes supplies the vanishing at \NSi{NS-F-P135-002}{\eta=0}.

During the first transition \NSi{NS-F-P135-003}{U=4\eta},
\NSi{NS-F-P135-004}{W=1-4L<0} and \NSi{NS-F-P135-005}{l\geq0}, so the angular source is at
least \NSi{NS-F-P135-006}{-h}. The finite transition therefore preserves a positive lower
bound for \NSi{NS-F-P135-007}{Q_s}, and \NSi{NS-F-P135-008}{b_s=0},
\NSi{NS-F-P135-009}{a\leq2}. In the axial transition, the angular source
\NSi{NS-F-P135-010}{S_q} is
\NSFormula{NS-F-P135-011}{display}{A.26}
\begin{equation}
 S_q=-h(1-2k\eta^2)+(D+dk)\eta J_0',\qquad
 Q_s\geq c(\eta^2+e^{-y})\quad(0\leq y\leq T_d).
 \tag{A.26}\label{eq:A.26}
\end{equation}
The second estimate follows from variation of constants, since
\NSi{NS-F-P135-012}{\eta J_0'\asymp\eta^2}, \NSi{NS-F-P135-013}{D+dk\geq D}, the value of
\NSi{NS-F-P135-014}{Q_s} at the inner endpoint is positive, and
\NSi{NS-F-P135-015}{h\ll e^{-T_d}}.

In this interval the geometric axial source is \NSi{NS-F-P135-016}{O(|\eta|)}. By
\eqref{eq:A.23}, its pressure source is \NSi{NS-F-P135-017}{O(E^2|\eta|)}: the unmodified
profile at larger radii decays at least as a constant times
\NSi{NS-F-P135-018}{e^{-\Delta y/2}} and satisfies
\NSi{NS-F-P135-019}{|\partial_\eta\log E|\leq|J_0'|}. The subsequent azimuthal-profile
correction may be omitted because it preserves the total pressure increment. Solving
\NSi{NS-F-P135-020}{N_s'+N_s=S_n} and using
\NSi{NS-F-P135-021}{E^2=P_1^2f^2e^{-y}} gives
\NSFormula{NS-F-P135-022}{display}{none}
\[
 \frac{|N_s|}{E^2}\leq C|\eta|
 \left(\frac{e^{T_d}}{P_*^2}+1+y\right)\leq C|\eta|(1+y).
\]
Now \NSi{NS-F-P135-023}{b_s=2k'\eta/E}, so \eqref{eq:A.26} gives
\NSi{NS-F-P135-024}{|b_sw|\leq Ce_d}. Also \NSi{NS-F-P135-025}{b_s^2\ll1} because
\NSi{NS-F-P135-026}{P_*>e^{T_d}}. As \NSi{NS-F-P135-027}{a=2}, choosing
\NSi{NS-F-P135-028}{M_d} first to make \NSi{NS-F-P135-029}{e_d} small and
\NSi{NS-F-P135-030}{P_*} afterward proves both strict inequalities in \eqref{eq:A.24}.

The intermediate power-law interval. During the axial transition, write
\NSi{NS-F-P135-031}{\mathcal A_X(U)=\bar k\eta}. Then
\NSi{NS-F-P135-032}{\bar k'=k-\bar k}, and the final eleven units with
\NSi{NS-F-P135-033}{k=0} give \NSi{NS-F-P135-034}{L\bar k<1/2}. Hence
\NSi{NS-F-P135-035}{W=1-L\bar k\geq1/2} on the next transition. Put
\NSi{NS-F-P135-036}{c_\eta=D\eta J_0'\asymp\eta^2}. With
\NSi{NS-F-P135-037}{U=0} the angular source on the transition is
\NSi{NS-F-P135-038}{(-l)W-h+c_\eta}. A positive lower bound for
\NSi{NS-F-P135-039}{Q_s}, independent of \NSi{NS-F-P135-040}{\lambda}, persists. The
bound on \NSi{NS-F-P135-041}{w} depends only on the parameters fixed before
\NSi{NS-F-P135-042}{\lambda}. Thus \NSi{NS-F-P135-043}{\sqrt\lambda|w|} is small and
\NSi{NS-F-P135-044}{b_s=0}.

In the local logarithmic radial coordinate on the constant-slope interval
\NSi{NS-F-P135-045}{\bar k(y)=\bar k(0)e^{-y}}, so the angular source is exactly
\NSi{NS-F-P135-046}{\lambda-h+c_\eta-\lambda L\bar k(0)e^{-y}}. Integrating the scalar
equation for \NSi{NS-F-P135-047}{Q_s} gives
\NSFormula{NS-F-P135-048}{display}{A.27}
\begin{equation}
\begin{aligned}
 Q_s&\geq c_{\mathrm{pre}}(\eta^2+\lambda+e^{-(1-\lambda)y}),\\
 Q_s-\frac{\lambda-h+c_\eta}{1-\lambda}
   &=O_{C^1_\eta}\!\left(C_{\mathrm{pre}}(1+y)e^{-(1-\lambda)y}\right),\\
 |N_s/E|&\leq C_{\mathrm{pre}}|\eta|(1+y)e^{-(1/2-\lambda)y}.
\end{aligned}
\tag{A.27}\label{eq:A.27}
\end{equation}
For the last line, \NSi{NS-F-P135-049}{U=0} eliminates the geometric source and the pressure
source remains \NSi{NS-F-P135-050}{O(E^2|\eta|)}; convolution with
\NSi{NS-F-P135-051}{e^{-y}} proves the bound. One \NSi{NS-F-P135-052}{\eta} derivative
obeys the same estimate without \NSi{NS-F-P135-053}{|\eta|}. Using
\NSi{NS-F-P135-054}{|\eta|/(\eta^2+e^{-(1-\lambda)y})\leq\tfrac12e^{(1-\lambda)y/2}}, we obtain
\NSFormula{NS-F-P135-055}{display}{none}
\[
 \sqrt\lambda|w|\leq C_{\mathrm{pre}}\sqrt\lambda(1+y)e^{\lambda y/2}=o(1)
 \qquad(0\leq y\leq60\log(1/\lambda)).
\]
Since \NSi{NS-F-P135-056}{a=2+2\lambda} and \NSi{NS-F-P135-057}{b_s=0}, this proves the
admissible stress cone there. On the preceding transition the same argument gives the
admissible stress cone condition whenever \NSi{NS-F-P135-058}{l<0}, and the relaxed
condition at its initial endpoint.

The pulse. Here both shear components vary, so the cone test requires a more precise relation
between the averaged axial velocity and the pulse profile. Set
\NSi{NS-F-P135-059}{m=\mathcal A_X(U)/E} and \NSi{NS-F-P135-060}{\beta=1/2-\lambda}. The
exact equation for \NSi{NS-F-P135-061}{m} is \NSi{NS-F-P135-062}{m'+\beta m=R_b}.
Extending the main pulse by zero on the left and expanding its exponential convolution twice
gives
\NSFormula{NS-F-P135-063}{display}{A.28}
\begin{equation}
 m=\frac{R_b}{\beta}-\frac{\lambda\partial_{\zeta_b}R_b}{\beta^2}
   +O(\lambda^2)+m(0)e^{-\beta y}.
 \tag{A.28}\label{eq:A.28}
\end{equation}
\NSPage{136}%
Indeed Taylor's formula bounds the convolution remainder by
\NSi{NS-F-P136-001}{C\lambda^2\lVert\partial_{\zeta_b}^2R_b\rVert_\infty
\int_0^\infty e^{-\beta v}v^2\,dv}. The start is flat, and the end bumps have exponentially
small fixed \NSi{NS-F-P136-002}{\zeta_b} derivatives. The same expansion holds after one
\NSi{NS-F-P136-003}{\eta} derivative. The initial value is bounded by \eqref{eq:A.14}.

In the angular equation, \NSi{NS-F-P136-004}{W=1-2D\eta Em-d(Em)_\eta} and
\eqref{eq:A.20} show that
\NSi{NS-F-P136-005}{S_q=\lambda-h+c_\eta+O(C_{\mathrm{pre}}E)}. Equations
\eqref{eq:A.14} and \eqref{eq:A.27} therefore imply
\NSFormula{NS-F-P136-006}{display}{A.29}
\begin{equation}
 Q_s=\frac{\lambda-h+c_\eta}{1-\lambda}+O(C_{\mathrm{pre}}\lambda^{28}).
 \tag{A.29}\label{eq:A.29}
\end{equation}
The pressure normalized to vanish at radial infinity and its first
\NSi{NS-F-P136-007}{\eta} derivative are \NSi{NS-F-P136-008}{O(E^2)}. Furthermore,
integrating \NSi{NS-F-P136-009}{XE^2(R_b^2-1/2)} and the moment discrepancy accumulated
before the pulse gives
\NSi{NS-F-P136-010}{S/(XE^2)=O_{C^1_\eta}(C_{\mathrm{pre}}e^{26}/\lambda)} throughout the
pulse. Formula~\eqref{eq:4.16} expresses \NSi{NS-F-P136-011}{Q_s,N_s} in terms of the radial
moment functions and their \NSi{NS-F-P136-012}{\eta}-derivatives, without radial derivatives
of \NSi{NS-F-P136-013}{E,U}. Substitution yields
\NSFormula{NS-F-P136-014}{display}{A.30}
\begin{equation}
 \frac{N_s}{E}=-R_b+(D+c_\eta)m-D\eta m_\eta+O(C_{\mathrm{pre}}\lambda^{27}).
 \tag{A.30}\label{eq:A.30}
\end{equation}
Here all constants are independent of \NSi{NS-F-P136-015}{\mathsf C}. Because
\NSi{NS-F-P136-016}{|\eta|/(\lambda+c_\eta)\leq C/\sqrt\lambda}, \eqref{eq:A.20} makes
the contribution of \NSi{NS-F-P136-017}{\eta m_\eta/Q_s} an
\NSi{NS-F-P136-018}{o(1)} term. Combining Equations \eqref{eq:A.28} to \eqref{eq:A.30}
gives
\NSFormula{NS-F-P136-019}{display}{none}
\[
 w=2R_b-C_d\partial_{\zeta_b}R_b+o(1),
 \qquad C_d=\frac{\lambda(D+c_\eta)(1-\lambda)}{\beta^2(\lambda-h+c_\eta)}\leq2+o(1).
\]
The estimate for \NSi{NS-F-P136-020}{w} now reduces the cone test to bounds on the pulse
shape and its derivative. For the coefficient of \NSi{NS-F-P136-021}{R_b} use
\NSi{NS-F-P136-022}{D+c_\eta-\beta=\lambda-h+c_\eta}; the bound on
\NSi{NS-F-P136-023}{C_d} follows by maximizing the quotient at
\NSi{NS-F-P136-024}{c_\eta=0}, with \NSi{NS-F-P136-025}{h/\lambda\to0}. Also
\NSi{NS-F-P136-026}{a=2+2\lambda} and
\NSi{NS-F-P136-027}{b_s=-R_b+O(\lambda)}. Apart from exponentially small end bumps,
\NSi{NS-F-P136-028}{R_b\geq0} and \NSi{NS-F-P136-029}{\partial_{\zeta_b}R_b\leq1.2}. Thus
\NSFormula{NS-F-P136-030}{display}{none}
\[
\begin{aligned}
 b_sw&\leq\sup_{R\geq0}(-2R^2+2.41R)+o(1)<.74,\\
 2b_sw+b_s^2/a+(a-2)w^2
   &\leq\sup_{R\geq0}(-3.5R^2+4.82R)+o(1)<1.68.
\end{aligned}
\]
Only the upper bound on the pulse derivative is used in these quadratic cross terms; its
negative derivative near the cutoff decreases the left-hand sides. All derivatives remain
bounded by fixed constants. The first margin gives \NSi{NS-F-P136-031}{a-b_sw>1}, and
\NSi{NS-F-P136-032}{v_s\geq a>2}, proving the admissible stress cone on the whole pulse.

Profile interpolation and exterior transition. After the pulse,
\NSi{NS-F-P136-033}{U=M=J=0} and hence \NSi{NS-F-P136-034}{W=1}. During the profile
interpolation the angular source is \NSi{NS-F-P136-035}{-l-h+\vartheta_f c_\eta}, so
\NSi{NS-F-P136-036}{Q_s\geq c\lambda}; the small \NSi{NS-F-P136-037}{C^1_\eta} angular
correction preserves this bound. Since \NSi{NS-F-P136-038}{S(\infty)=0},
\NSFormula{NS-F-P136-039}{display}{A.31}
\begin{equation}
 \frac{S(X)}{X}=\frac1{2X}\int_X^\infty E(X')^2\,dX'.
 \tag{A.31}\label{eq:A.31}
\end{equation}
Until the exterior transition, \NSi{NS-F-P136-040}{l\leq-\lambda/2}, and the remaining
exterior transition integral has the bound \eqref{eq:A.18}. Thus
\NSi{NS-F-P136-041}{(|S|+|S_\eta|)/X\leq CE^2/\lambda}, including the angular edit, while
\NSi{NS-F-P136-042}{|\Pi|+|\Pi_\eta|\leq CE^2}. Formula~\eqref{eq:4.16} now gives
\NSi{NS-F-P136-043}{|w|\leq CE/\lambda^2\ll1}, since pulse end made
\NSi{NS-F-P136-044}{E} exponentially small in \NSi{NS-F-P136-045}{1/\lambda}. Here
\NSi{NS-F-P136-046}{b_s=0} and
\NSi{NS-F-P136-047}{2<a\leq2+2\lambda+.2+o(1)}, so both cone margins persist.

Once the exterior transition has started, its fields and remaining radial moment functions are
uniform in \NSi{NS-F-P136-048}{\eta} except for the explicit coefficients in the identity for
\NSi{NS-F-P136-049}{Q_s,N_s}. Combining
\NSPage{137}%
\eqref{eq:A.23} and \eqref{eq:A.31} gives
\NSFormula{NS-F-P137-001}{display}{none}
\[
 N_s=2\eta\int_y^\infty\bigl(he^{v-y}-A\bigr)E(v)^2\,dv=O(E(y)^2).
\]
Throughout the exterior transition and at larger radii,
\NSi{NS-F-P137-002}{l\leq-3h/4}. Since
\NSi{NS-F-P137-003}{(\log E)'=l-1/2}, for \NSi{NS-F-P137-004}{v\geq y} we have
\NSFormula{NS-F-P137-005}{display}{none}
\[
 E(v)^2=E(y)^2\exp\!\left(\int_y^v(2l(s)-1)\,ds\right)
 \leq E(y)^2e^{-(1+3h/2)(v-y)}.
\]
Consequently
\NSFormula{NS-F-P137-006}{display}{none}
\[
\begin{aligned}
 h\int_y^\infty e^{v-y}E(v)^2\,dv
   &\leq hE(y)^2\int_0^\infty e^{-3hs/2}\,ds=\frac23E(y)^2,\\
 \int_y^\infty E(v)^2\,dv
   &\leq E(y)^2\int_0^\infty e^{-(1+3h/2)s}\,ds
     =\frac{E(y)^2}{1+3h/2}.
\end{aligned}
\]
With \NSi{NS-F-P137-007}{A=1/2+h} and \NSi{NS-F-P137-008}{|\eta|\leq1}, these estimates
give the stated bound for \NSi{NS-F-P137-009}{N_s} uniformly in
\NSi{NS-F-P137-010}{h}. Before the steep-decay interval
\NSi{NS-F-P137-011}{Q_s\geq c\lambda}. After it, and through the first half-unit of the
terminal interval, \NSi{NS-F-P137-012}{Q_s\geq ch}: in the last constant-slope interval it
decreases only to \NSi{NS-F-P137-013}{Q_p\asymp h}, and \eqref{eq:A.16} gives the same
lower bound for \NSi{NS-F-P137-014}{0\leq y\leq1/2}. By \eqref{eq:A.17},
\NSi{NS-F-P137-015}{E\leq Ce_{\mathrm{rel}}h^6} after the steep-decay interval, so
\NSi{NS-F-P137-016}{E/Q_s\ll1} even when \NSi{NS-F-P137-017}{h} is much smaller than
\NSi{NS-F-P137-018}{e_{\mathrm{rel}}}. Consequently \NSi{NS-F-P137-019}{|w|\ll1},
\NSi{NS-F-P137-020}{b_s=0}, and \NSi{NS-F-P137-021}{2<a\leq4} on the exterior transition
interval under consideration. This proves the strict true cone there.

Every lower bound for \NSi{NS-F-P137-022}{Q_s} used here is on a fixed compact logarithmic
interval after all choices in \eqref{eq:A.6}. The fields and the functions
\NSi{NS-F-P137-023}{Q_s,N_s} on such an interval are independent of
\NSi{NS-F-P137-024}{X_R} under the normalizations \eqref{eq:A.4}. Hence one final increase
of \NSi{NS-F-P137-025}{X_R} makes the large-\NSi{NS-F-P137-026}{p_{s,1}} test uniform.
This finishes the proof of Proposition~\ref{prop:A.4}.

\subsection{Replacing the exterior power by a heat flow.}\label{sec:A.6}
The reference profile \NSi{NS-F-P137-027}{(U,E)} of Proposition~\ref{prop:A.4} has the
required moments and satisfies \NSi{NS-F-P137-028}{U=0},
\NSi{NS-F-P137-029}{E=E_{\mathrm{pow}}=c_\infty X^{-A}} beyond the terminal interval. The
exterior power law \NSi{NS-F-P137-030}{E_{\mathrm{pow}}} still has a viscous residual. In
this subsection we will replace it beyond the annulus by the exact radial swirl heat flow
\NSi{NS-F-P137-031}{K} of Lemma~\ref{lem:A.6}. This will make the exterior momentum
residual vanish and give the field smooth limits at the singular time away from the axis. This
replacement changes three moments; a correction in the second reserved interval
\NSi{NS-F-P137-032}{\mathcal I_2} restores them (Proposition~\ref{prop:A.7}), including the pressure
datum \NSi{NS-F-P137-033}{\Pi_0} already prescribed for the axis in \eqref{eq:A.21}.

For a profile regular at the axis, write
\NSi{NS-F-P137-034}{\mathscr R_\theta^{(0)},\mathscr R_z^{(0)}} for its leading azimuthal
and axial momentum residuals. Thus \NSi{NS-F-P137-035}{\mathscr R_j^{(0)}=R_j^{(0)}} in
the notation of \eqref{eq:4.12}; these include radial viscosity and omit the higher-order axial
viscosity. The stress components are
\NSi{NS-F-P137-036}{T_\theta(r)=-r^{-2}\int_0^r s^2\mathscr R_\theta^{(0)}(s)\,ds} and
\NSi{NS-F-P137-037}{T_z(r)=-r^{-1}\int_0^r s\mathscr R_z^{(0)}(s)\,ds}, as in
Proposition~\ref{prop:4.2}. In the following subsection we will use the corrected moment
conditions to rewrite these integrals as integrals from \NSi{NS-F-P137-038}{r} to infinity
(Lemma~\ref{lem:A.8}) and determine the sign and limiting direction of the stress at the outer
endpoint (Proposition~\ref{prop:A.10}).

Throughout this subsection \NSi{NS-F-P137-039}{0<h<1/2},
\NSi{NS-F-P137-040}{A=1/2+h}, and \NSi{NS-F-P137-041}{D=1/2-h}. We retain
\NSi{NS-F-P137-042}{d=1-\eta^2}, \NSi{NS-F-P137-043}{\tau=1-t}, and
\NSi{NS-F-P137-044}{L=1-2h\eta^2} from \eqref{eq:4.1}. All parameters in the construction
of the outer radial profile are fixed in the order \eqref{eq:A.6} before
\NSi{NS-F-P137-045}{X_R} is increased. In the local logarithmic radial coordinate for the
tail of \eqref{eq:A.12}, write \NSi{NS-F-P137-046}{y=\log(X/X_{\mathrm{tail}})},
\NSi{NS-F-P137-047}{X_b=e^3X_{\mathrm{tail}}}, and
\NSi{NS-F-P137-048}{E_{\mathrm{pow}}=c_\infty X^{-A}}, where
\NSi{NS-F-P137-049}{c_\infty>0} is independent of \NSi{NS-F-P137-050}{\eta}. Thus the
reference tail is \NSi{NS-F-P137-051}{E_{\mathrm{pow}}f_o(y)}.

\NSPage{138}%
Set \NSi{NS-F-P138-001}{a_K=1+h}, and define the heat factor by
\NSFormula{NS-F-P138-002}{display}{A.32}
\begin{equation}
 \mathcal H(Z)=\frac1{\Gamma(a_K)}\int_0^\infty
 e^{-v}v^{a_K-1}(1+Zv)^{-h}\,dv,\qquad Z\geq0,
 \tag{A.32}\label{eq:A.32}
\end{equation}
Here \NSi{NS-F-P138-003}{\Gamma(a_K)=\int_0^\infty e^{-v}v^{a_K-1}\,dv} is the gamma
integral, so \NSi{NS-F-P138-004}{\mathcal H(0)=1}. Write
\NSi{NS-F-P138-005}{E_{\mathrm{heat}}(X,\eta)=E_{\mathrm{pow}}(X)\mathcal H(2d/X)} for the
heat profile.

\begin{lemma}\label{lem:A.6}
The function \NSi{NS-F-P138-006}{\mathcal H} is positive and smooth up to
\NSi{NS-F-P138-007}{Z=0} from the right. Every fixed derivative is bounded on bounded
nonnegative intervals. The field
\NSFormula{NS-F-P138-008}{display}{A.33}
\begin{equation}
 K(r,t)=c_\infty s^{-A}\mathcal H(2\tau/s),\qquad s=r^2/2,
 \tag{A.33}\label{eq:A.33}
\end{equation}
is independent of \NSi{NS-F-P138-009}{z}, satisfies
\NSi{NS-F-P138-010}{\partial_tK=(\partial_{rr}+r^{-1}\partial_r-r^{-2})K}, and has
\NSi{NS-F-P138-011}{K_r<0}. Its profile is
\NSi{NS-F-P138-012}{E_{\mathrm{pow}}(X)\mathcal H(2d/X)}, smooth with all
\NSi{NS-F-P138-013}{\eta}-derivatives extending continuously to
\NSi{NS-F-P138-014}{\eta=\pm1} from the interior.
\end{lemma}
\begin{proof}
For the rising factorial
\NSi{NS-F-P138-015}{(b)_m=b(b+1)\cdots(b+m-1)}, with
\NSi{NS-F-P138-016}{(b)_0=1}, dominated differentiation gives
\NSFormula{NS-F-P138-017}{display}{A.34}
\begin{equation}
 \mathcal H^{(m)}(Z)=\frac{(-1)^m(h)_m}{\Gamma(1+h)}
 \int_0^\infty e^{-v}v^{h+m}(1+Zv)^{-h-m}\,dv,
 \tag{A.34}\label{eq:A.34}
\end{equation}
\NSFormula{NS-F-P138-018}{display}{A.35}
\begin{equation}
 \mathcal H^{(m)}(0)=(-1)^m(h)_m(1+h)_m.
 \tag{A.35}\label{eq:A.35}
\end{equation}
The integrable majorant obtained by deleting the last factor works also at
\NSi{NS-F-P138-019}{Z=0}. In particular, for each fixed derivative order the usual finite
Taylor formula is valid there; no convergent Taylor series is needed. On a fixed small
nonnegative interval it gives
\NSFormula{NS-F-P138-020}{display}{A.36}
\begin{equation}
 \mathcal H(Z)=1-h(1+h)Z+O_h(Z^2),\qquad
 |\mathcal H(Z)-1|+|Z\mathcal H'(Z)|\leq C_hZ.
 \tag{A.36}\label{eq:A.36}
\end{equation}
Here the last constant can be uniform for \NSi{NS-F-P138-021}{0<h\leq h_0<1/2}. Indeed
\NSi{NS-F-P138-022}{(h)_m(1+h)_m\leq C_mh} for every fixed
\NSi{NS-F-P138-023}{m\geq1}. Integration of
\NSi{NS-F-P138-024}{h\partial_v[e^{-v}v^{a_K}(1+Zv)^{-a_K}]} gives
\NSFormula{NS-F-P138-025}{display}{A.37}
\begin{equation}
 Z^2\mathcal H''+(1+2a_KZ)\mathcal H'+a_K(a_K-1)\mathcal H=0;
 \tag{A.37}\label{eq:A.37}
\end{equation}
both boundary terms vanish. Direct differentiation of \eqref{eq:A.33}, with
\NSi{NS-F-P138-026}{Z=2\tau/s}, gives
\NSFormula{NS-F-P138-027}{display}{none}
\[
\begin{aligned}
 \partial_tK&=-2c_\infty s^{-A-1}\mathcal H',\\
 (\partial_{rr}+r^{-1}\partial_r-r^{-2})K
 &=2c_\infty s^{-A-1}
   [Z^2\mathcal H''+(2A+1)Z\mathcal H'+(A^2-1/4)\mathcal H].
\end{aligned}
\]
Since \NSi{NS-F-P138-028}{2A+1=2a_K} and
\NSi{NS-F-P138-029}{A^2-1/4=a_K(a_K-1)}, \eqref{eq:A.37} proves the heat equation with its
swirl term. Moreover, \NSi{NS-F-P138-030}{-Z\mathcal H'/\mathcal H} is the average of
\NSi{NS-F-P138-031}{hZv/(1+Zv)} against the positive probability density proportional to the
integrand in \eqref{eq:A.32}. Consequently
\NSFormula{NS-F-P138-032}{display}{none}
\[
 0\leq-\frac{Z\mathcal H'}{\mathcal H}<h,
 \qquad
 \frac{rK_r}{2K}=-A-\frac{Z\mathcal H'}{\mathcal H}<-\frac12.
\]
The composition with \NSi{NS-F-P138-033}{Z=2(1-\eta^2)/X} uses only
\NSi{NS-F-P138-034}{Z\geq0}. Each nonzero \NSi{NS-F-P138-035}{\eta}-derivative is a finite
sum of terms \NSi{NS-F-P138-036}{\mathcal H^{(k)}(2d/X)(2/X)^k} times polynomials in
\NSi{NS-F-P138-037}{d'=-2\eta} and \NSi{NS-F-P138-038}{d''=-2}. Thus no division by
\NSi{NS-F-P138-039}{d} occurs. This proves the claimed closed-interval smoothness. The
integral formula is used only on its stated nonnegative domain. In particular the large-
\NSi{NS-F-P138-040}{X} expansion, with its remainder after every fixed application of
\NSi{NS-F-P138-041}{((X\partial_X)^j\partial_\eta^m)}, is
\NSFormula{NS-F-P138-042}{display}{A.38}
\begin{equation}
 \frac{E_{\mathrm{heat}}(X,\eta)}{E_{\mathrm{pow}}(X)}
 =\mathcal H(2d/X)=1-\frac{2h(1+h)d}{X}+O_{h,j,m}(X^{-2}).
 \tag{A.38}\label{eq:A.38}
\end{equation}
\end{proof}
